AI-assisted animation systems increasingly support generation and low-level refinement, yet conversational character performance also depends on interpretive decisions about what each character is doing, attending to, and expressing over time. We present a context-aware authoring workflow that exposes these decisions as editable \textit{Semantic Performance Tags} between narrative interpretation and animation execution. For dyadic dialogue, a pretrained large language model proposes an initial performance state for both characters and sparse dialogue-anchored semantic beats; the Maya interface projects the resulting canonical plan into two character-specific tag panels where animators can inspect and revise acting-related affect, attention, head, and blink decisions before deterministic compilation to JALI and Maya controls. The representation supports revising one or multiple semantic decisions across either character while keeping the dialogue immutable and avoiding an additional LLM call for tag edits. We evaluate the approach in two complementary studies: a within-subject authoring study comparing Direct Generation with Editable Semantic Performance Tags, and a text-only context ablation examining how additional conversational context affects proposed performance decisions. Results show \textbf{[headline Study~A / RQ2--RQ3 result]}, while the context ablation finds \textbf{[headline Study~B / RQ1 result]}. These findings suggest \textbf{[replace with a results-supported implication after analysis]}.
